% 共用表格欄位；由摘要及第一章載入，避免修改主檔。
\ifdefined\proposalsetupdone\else
\gdef\proposalsetupdone{1}
\newcolumntype{P}[1]{>{\raggedright\arraybackslash\hspace{0pt}}p{#1}}
% 書目依使用者指定：中文文獻在前，英文文獻在後。
\nocite{taichung_eat_outside,taichung_gender,tdx_planner,tdx_bus,llm_planning_tools,subjective_attributes,critique_preference,bus_dynamic_replanning,llm_planning,ttg,hypertree_planning,bright,hybrid_retrieval,knowledge_graph_llm_recommendation,restaurant_habib,poi_stisan,bus_delay_uncertainty,google_places_policy}
\fi

\chapter{Introduction 緒論}

\section{研究背景與動機}

外食是日常生活中的重要選擇。臺中市政府於 2025 年發布的家庭外食分析指出，2024 年底臺中市餐飲業營利事業家數達 22,612 家，較 2015 年增加 47.55\%。面對多樣的選擇，使用者在約會、慶生或朋友聚餐時，需要知道哪些餐廳符合當次用餐目的~\cite{taichung_eat_outside}。

餐廳評論常包含「適合聊天」「氣氛浪漫」「慶生有儀式感」等情境描述。這些語句與使用者需求未必使用相同詞彙，適合透過語意方法找出相關證據。主觀屬性與餐廳對話研究，也說明自然語言偏好與評論片段對推薦的價值~\cite{subjective_attributes,critique_preference}。

使用者還需要實際前往餐廳。本研究以臺中市的公車用餐情境為例，考慮從出發地到候選餐廳的步行、候車、乘車與轉乘。臺中市政府於 2026 年發布的運具通報，以 2024 年調查資料指出，男性與女性使用汽車客運的市占率分別為 3.9\% 與 6.0\%。該類別涵蓋多種客運服務，不直接等同市區公車比例；本研究以公車作為交通分析對象~\cite{taichung_gender}。

在技術面，自然語言規劃研究已將需求轉換與形式化查核分工，公共運輸研究也以班次、轉乘與步行模型處理旅程。這些方法使系統能在找到語意相關的餐廳之後，進一步比較交通條件，將交通便利性回饋到餐廳的最終推薦~\cite{llm_planning_tools,bus_dynamic_replanning}。

\section{問題陳述與既有方法的不足}

以「今晚要和朋友慶生，每人 600 元內，從台中火車站搭公車去」為例，系統首先要理解慶生的情境，從餐廳評論找到有相關依據的候選；接著，才能依使用者的位置與時間安排公車交通，排除不符合條件或交通負擔過大的選擇。

本研究要解決的問題是：如何在保留情境相關性的前提下，利用臺中公車交通條件篩選候選餐廳，使最終推薦同時符合「適合這次用餐」與「方便前往」？當店家或交通資料不足時，又應如何呈現未知項，避免推薦文字承諾超出證據？

只做語意排名，不能保證交通方便；只找交通最近的餐廳，也不能回答約會、慶生等情境需求。既有研究分別提供評論偏好與交通分析方法，本研究將兩者接合成先召回、再交通篩選的流程，並評估交通分析是否實際改變推薦品質~\cite{critique_preference,bus_dynamic_replanning}。

\section{本研究的方法與創新}

為此，本研究提出一個考量臺中公車交通便利性的情境感知餐廳推薦 AI Agent System。語言模型先將自然語言需求整理成情境、預算、出發地、時間及交通限制；檢索模組將情境與餐廳評論語意對齊，召回一批有來源支持的餐廳候選。

對每間候選餐廳，系統再呼叫交通工具取得公車方案，由程式計算步行、候車、乘車與轉乘負擔，檢查抵達後是否仍可完成用餐。明確不符合必要條件者排除；資料不足者標為不確定；通過者再結合情境相關性與交通便利性排序。

本研究的設計差異是讓交通分析在最終選餐廳之前發生，並以可追溯的判定紀錄說明推薦取捨。向量檢索、知識圖譜與公車規劃均有既有方法，本研究不將其個別視為新演算法，而著重於流程整合與本場域的評估。

\section{研究問題}
\label{sec:research_questions}

\begin{itemize}
\item RQ1：向量或混合檢索是否比關鍵字，更能將約會、慶生等情境與餐廳評論對齊？
\item RQ2：在候選形成後加入交通篩選與排序，是否比只依語意排名、最後才補路線，更能兼顧用餐情境與交通便利性？
\item RQ3：三值查核與來源、時間紀錄，是否能減少違反必要條件及證據不足卻宣稱符合的情形？
\item RQ4：完整時間線是否比只看出發時營業或 API 的單一時間，更能避免無法上車、轉乘不成立與用餐超過打烊？
\end{itemize}

上述問題界定系統設計與方法分析的重點。

\section{預期研究貢獻}

\begin{enumerate}
\item 建立將用餐情境與餐廳評論語意對齊的候選召回流程，推薦時可回到具體評論依據。
\item 在候選餐廳的最終選擇前加入臺中公車交通分析，比較步行、候車、乘車與轉乘，讓交通便利性影響餐廳排序。
\item 以三值查核及完整時間線處理預算、營業與搭乘條件，區分不符合與證據不足。
\item 提供可回溯評論與交通資料的推薦說明，協助使用者理解餐廳選擇的依據與交通取捨。
\end{enumerate}

\section{論文架構}

第二章回顧 AI 技術與餐廳、公車領域研究，整理文獻彙整表與研究缺口。第三章歸納餐廳與公共運輸推薦的核心設計觀點，討論關鍵方法取捨，並以推理問題釐清領域知識。第四章提出系統架構、資料與時間模型、判定與排序、核心演算法、流程，最後整理預期成果與後續擴充。
